\begin{textsample}{POS Dim 1 – vem\_ed\_32 – Score 38.00 – vem\_ed\_32/t170.txt}  \label{ex:f1_pos_032}
Mesa-redonda virtual debate benefícios dos projetos COIL Maria Claudia Nunes Delfino e Neusa Haruka Sezaki Gritti, coordenadoras de Projetos Colaborativos Internacionais (PCIs) em inglês na Área de Apoio à Internacionalização do Ensino Superior da Divisão de Extensão e Pesquisa (Depes) da Coordenadoria de Ensino Superior de Graduação (CGESG) do Centro Paula Souza (CPS), e Tálita Guarino, professora de inglês da Fatec Guaratinguetá, participaram de mesa-redonda virtual no XIV Congresso Internacional de Ciência, Tecnologia e Desenvolvimento (CICTED), evento gratuito \textbf{promovido} pela Universidade de Taubaté (Unitau) em outubro de 2025.
A mesa-redonda \textbf{foi} composta ainda por Anoush Ayunts, da Yerevan State University (Armênia) e \textbf{contou} com a moderação de Adriana Leônidas de Oliveira (Unitau).
Neusa trouxe as definições de Intercâmbio Virtual e da abordagem COIL.
Claudia resumiu as fases dos projetos COIL e em sua articulação para a busca de solu- ções para problemas \textbf{locais} ou globais. “COIL se adapta a múltiplas \textbf{disciplinas}: desde administração até estudos \textbf{ambientais} e ciências da saúde, mas o \textbf{que} une \textbf{todas} essas experiências \textbf{é} o \textbf{foco} na colaboração e \textbf{compreensão} intercultural”.
As palestrantes apresentaram os benefícios culturais e \textbf{linguísticos}, para estu- dantes de graduação, dos projetos COIL – \textbf{que}, nas Fatecs do CPS, recebem o nome de PCIs.
A professora Anoush realiza, desde 2021, um PCI entre a Yerevan State University e Fatec São José do Rio Preto, com a parceria da professora Edilene Gasparini Fernandes, \textbf{que} incluiu Webquests no projeto (leia Artigo de Opinião em VEm 25: https://tinyurl.com/5d9pfp6k).
De 2022 em diante, o projeto passou a abordar Objetivos de Desenvolvimento \textbf{Sustentável} da ONU e se es- tendeu para as unidades de Bragança Paulista, Garça, Jales, Mogi das Cruzes, Pompeia e Guaratinguetá (com Tálita Guarino).
Anoush comentou sobre uma publicação derivada dessa experiência, na Revista CBTecLE, o depoimento do aluno Israel Nascimento para o canal da UNICollaboration no YouTube e a publicação no site do CPS sobre o PCI.
Tálita Guarino \textbf{compartilhou} depoimentos em vídeo de três estudantes e elencou benefícios do PCI para alunos, professores e instituições de ensino \textbf{superior} (leia mais sobre o assunto no Artigo de Opinião).
A professora \textbf{identificou} quatro be- nefícios principais a \textbf{partir} dos relatos dos discentes: Colaboração intercultural – \textbf{aprender} a respeitar e valorizar diferentes \textbf{perspectivas}; \textbf{Desenvolvimento} \textbf{linguístico} - usar o inglês como ferramenta para \textbf{conexão}, não para perfeição; \textbf{Competência} digital – dominar ferramentas tecnológicas; Aprendizagem \textbf{socioemocional} – empatia, tolerância, \textbf{flexibilidade} e confiança. “Os estudantes mudaram seu mindset – pararam de se preocupar com erros e \textbf{começaram} a focar em \textbf{interação} \textbf{significativa}”, comentou Tálita.
Ela elencou ainda os benefícios para os professores, em uma \textbf{transformação} pedagógica: de “entregadores de \textbf{conteúdo}” para facilitadores de aprendizagem.
Isso envolve um planejamento intencional junto com o parceiro da IES internacional, gerenciamento das \textbf{diferenças} culturais e de fusos \textbf{horários}.
Para as instituições, o COIL \textbf{promove} Internacionalização em Casa de \textbf{forma} autêntica.

% matched lemmas: ambiental, aprender, começar, compartilhar, competência, compreensão, conexão, contar, conteúdo, desenvolvimento, diferença, disciplina, flexibilidade, foco, forma, horário, identificar, interação, linguístico, local, partir, perspectiva, promover, que, ser, significativo, socioemocional, superior, sustentável, todo, transformação
\end{textsample}
